\documentclass[11pt]{article}
\usepackage{amsmath,amssymb,amsthm,mathtools}
\usepackage[hidelinks]{hyperref}

\newtheorem{lemma}{Lemma}
\newtheorem{theorem}{Theorem}

\title{Online Dynamic Resource Allocation with State Constraints}
\author{Kamiar Asgari \and Michael J. Neely}
\date{April 11, 2026}

\begin{document}
\maketitle

\begin{abstract}
We study online dynamic resource allocation with resource-in-storage availability constraints under adversarial inflows and convex costs through the lens of adversarial online control. At each round, the controller chooses an allocation within available resources before observing the inflow and cost. We introduce \emph{Simplex Disturbance-Action Control} (SDAC), a simplex-constrained disturbance-action policy class that preserves feasibility by construction, and we develop an online controller over SDAC weights. We benchmark this algorithm against the best fixed SDAC policy in hindsight. Our regret analysis yields sublinear dependence on the horizon and logarithmic dependence on the policy dimension. The setup is motivated by applications in energy-harvesting battery management where decisions must satisfy storage-availability limits in real time.
\end{abstract}

\section{Introduction}
\label{sec:introduction}

\subsection{Motivation and Problem Setting}
Many dynamic resource systems are governed by a budget state variable that evolves over time and constrains the allocation action at every round. The controller must choose an allocation action from currently available resources before observing the round's inflow and cost.

This creates an online optimization problem with state coupling: aggressive early allocation can harm later performance, while overly conservative allocation can increase immediate cost. Our objective is to design a policy that remains feasible at every round and competes with the best fixed policy in a structurally tractable comparator class.

\subsection{Model and Information Structure}
We study a discrete-time storage-allocation system with state $b_t \geq 0$, where $b_t$ represents the currently available resource. At each time $t$, the controller chooses an allocation action $u_t$ under the \emph{feasibility constraints} such as
\[
0 \leq u_t \leq b_t,
\]
which means the controller cannot allocate more than the currently available resource or any negative amount. After the action is chosen, an adversarial inflow $a_t$ arrives, and the state evolves according to
\[
b_{t+1} = b_t - u_t + a_t,
\]
with inflow $a_t$.

Unlike stochastic application models that rely on distributional assumptions, our formulation is adversarial. Inflows $a_t$ and convex costs $c_t(b, u)$ may vary arbitrarily over time, and the controller chooses $u_t$ before observing either object at round $t$. The goal is low cumulative cost relative to the best fixed policy in a later-defined comparator class, while maintaining feasibility throughout.

\subsection{Application Domains}
This abstraction captures a broad set of domains:
\begin{itemize}
    \item \textbf{Energy Harvesting and Battery Management:} $b_t$ is battery state of charge, $a_t$ is harvested energy, and $u_t$ is discharge or consumption.
    \item \textbf{Queueing Networks and Buffer Management:} $b_t$ is queue/buffer occupancy, $a_t$ is arrivals, and $u_t$ is service rate constrained by current backlog.
    \item \textbf{Inventory Control and Supply Chains:} $b_t$ is on-hand inventory, $a_t$ is replenishment, and $u_t$ is allocation/sales under stock constraints.
    \item \textbf{Investments and Bank Accounts:} $b_t$ is liquid capital/balance, $a_t$ is incoming cash flow, and $u_t$ is capital deployment or withdrawal under budget limits.
    \item \textbf{Water Reservoir and Dam Management:} $b_t$ is stored water, $a_t$ is inflow, and $u_t$ is controlled release under storage limits.
\end{itemize}


\subsection{Related Work}
\paragraph{Queueing and Network Optimization}
The queueing-network literature has extensively studied systems with arrivals, storage (queues), and controlled service/allocations. A classical line of work focuses on \emph{adversarial queueing theory} (AQT), where arrivals and routing decisions may be chosen adversarially and the primary objective is to establish stability (i.e., bounded queue lengths) under appropriate load conditions \cite{borodin2001adversarial, gamarnik1998stability, kiwi2005survey}. While this model shares structural similarities with our setting in allowing worst-case inputs, its scope is fundamentally different: the goal is stability rather than optimization of a cost function, and no notion of regret or comparison to an optimal dynamic policy is considered. Moreover, in many classical adversarial queueing formulations, actions are \emph{discrete} (e.g., routing or scheduling decisions), whereas our setting considers \emph{continuous allocation decisions}.

A second major line is \emph{stochastic network optimization} based on Lyapunov-drift arguments \cite{neely2010stochastic, georgiadis2006resource}. These works assume stochastic arrivals and network conditions and design policies (e.g., max-weight or drift-plus-penalty) that achieve near-optimal time-average utility while ensuring queue stability. In these formulations, the stage utility (equivalently, negative cost) is modeled primarily as a function of the current control action, with queue states entering through stability constraints and drift terms rather than through a joint state-action adversarial loss. In addition, many formulations use discrete or finitely parameterized action sets (e.g., selecting from a finite set of rates or schedules), in contrast to our continuous decision space.

A related direction is online convex optimization with long-term constraints using virtual queues to handle long-term constraints \cite{mahdavi2012trading, jenatton2016adaptive, yu2017online, JMLR:v21:16-494}. These works allow adversarial convex costs and provide regret and long-term constraint violation guarantees. However, the queueing structure is typically auxiliary, serving to enforce constraints rather than modeling a system's dynamic, and the objective does not explicitly depend on the evolving queue state. Furthermore, the benchmarks are standard regret notions relative to fixed comparators, rather than policies that account for state evolution.




\paragraph{Queueing Control / Buffer Management}

An alternative but mathematically parallel perspective comes from queueing control and buffer management in computer networks and telecommunications \cite{kleinrock1975queueing,leboudec2001network}. In that view, one tracks a finite-capacity queue with stochastic arrivals and real-time service decisions. A standard queue update is
\[
Q_{t+1} \,=\, \min\!\bigl\{K,\,\max\{0,\,Q_t-\mu_t\}+A_t\bigr\},
\]
where \(Q_t\) is queue state, \(A_t\) is arrival, \(\mu_t\) is service, and \(K\) is capacity. Under the feasibility condition \(0\le \mu_t\le Q_t\), this reduces to
\begin{equation*}
Q_{t+1} \;=\; \min\{\,K,\; Q_t-\mu_t+A_t\,\},
\end{equation*}
which is isomorphic to \eqref{eq:battery-evo} via \(Q_t\leftrightarrow b_t\), \(A_t\leftrightarrow a_t\), \(\mu_t\leftrightarrow u_t\), and \(K\leftrightarrow B\).

A representative cost function in this setting is
\[
c_t(q,s) \;=\; q \;+\; f(s),
\]
which penalizes both large backlog (queue length) and service effort, where \(f\) is a convex “penalty” function of the service rate.  
The objective is to minimize the long-run average of these costs, i.e., the sum of the average queue length and the average service penalty. This queueing interpretation is also the basis of Lyapunov drift methods for real-time constrained control \cite{neely2010stochastic}, and has been explicitly used in energy-harvesting communication models that treat battery state as an energy queue \cite{sudevalayam2011energy,sharma2010optimal,yang2011optimal}.



\paragraph{Energy Harvesting and Battery Management}
Energy-harvesting control has been studied primarily under stochastic assumptions and utility-maximization objectives, rather than worst-case regret. Prior work formulates an online convex optimization model with finite battery capacity and i.i.d. energy arrivals, showing $O(\sqrt{T})$ regret with a $O(\sqrt{T})$-sized battery \cite{Asgari2020}. Related studies consider a single-device energy harvesting system with stochastic channel and energy states, and combine drift-plus-penalty with online gradient updates to achieve utility within $O(\varepsilon)$ of optimum using $O(1/\varepsilon)$ battery size \cite{Yu2019}. Complementary survey work synthesizes energy-harvesting node architectures, ambient energy sources, storage technologies, and system-level design implications under intermittent recharge opportunities \cite{sudevalayam2011energy}. Stochastic-control formulations for harvesting sensor nodes also characterize optimal energy-management policies that stabilize operation and optimize long-run communication performance under energy-availability constraints \cite{sharma2010optimal}. In addition, offline scheduling results for energy-harvesting links develop structure-exploiting optimal packet-transmission policies under energy-causality and data-availability constraints \cite{yang2011optimal}. Related work on stochastic energy-harvesting networks analyzes long-run utility/throughput under Lyapunov-style scheduling with energy-availability constraints \cite{Huang2013}, while fixed-fraction style policies are shown to be near-optimal for single-user energy harvesting systems with i.i.d. harvests and concave utilities \cite{arafa2017energy}.

In contrast, our framework allows fully adversarial inflow and cost sequences and supports general convex costs that depend jointly on action and stored resource state. Therefore, unlike prior energy harvesting online formulations that rely on stochasticity, our setting targets worst-case performance guarantees.






\paragraph{Energy Storage and Energy Management Systems}
A distinct and highly active formulation is found in the modern study of \emph{Energy Storage and Microgrid Management}, where the abstract reservoir is physically realized as a battery's State of Charge (SoC). The system is driven by stochastic energy arrivals from renewable sources (solar, wind) and governed by strict electrochemical constraints: finite capacity, maximum charge/discharge power limits, and strict non-negativity. Researchers have synthesized classical storage theories with power systems engineering to leverage Markov Decision Processes, Lyapunov optimization, and Model Predictive Control for optimal energy dispatch \cite{sudevalayam2011energy,urgaonkar2011optimal,ozel2011optimal}. The focus is on balancing immediate economic utility of allocating power against opportunity costs of future scarcity, all while satisfying bounded state constraints. Recent work extends these techniques to networked microgrids with spatially distributed renewable sources and storage \cite{sun2014real,yang2011optimal}. While our online learning formulation addresses the adversarial uncertainty beyond the stochastic models typically assumed in modern energy systems, the structural parallels---especially the prominence of threshold-based policies, Lyapunov drift arguments, and the centrality of capacity constraints---demonstrate how classical optimization insights resurface across physical instantiations of the dynamic allocation problem.


\paragraph{Inventory Management}
Yang et al. study an online inventory problem with time-varying linear prices and demand constraints, establishing optimal competitive-ratio guarantees under linear costs \cite{yang2020online}, while Hihat et al. formulate multi-product inventory control with stateful dynamics and general convex losses under a mild demand non-degeneracy condition \cite{hihat2023inventory}. Other online-learning variants include mirror-descent methods for multiproduct lost-sales systems with censored feedback \cite{guo2024omd} and adversarial repeated newsvendor learning with censored demand \cite{lugosi2024hardness}. In parallel, the classical operations-research inventory-control literature develops structurally similar state dynamics, e.g., $I_{t+1}=\max(0,I_t+x_t-d_t)$, and characterizes optimal threshold policies through dynamic programming \cite{arrow1958studies,scarf1960optimality,zipkin2000foundations,porteus2002foundations}. From the competitive-analysis and online-optimization side, storage and inventory-constrained formulations are studied under worst-case models and linear objectives known at decision time \cite{Chau2016,Yang2020,yang2020online,lin2019competitive,chau2016cost,neely-inventory-control-cdc2010}, with related ideas in one-way trading and online knapsack \cite{cao2020optimal,el2001optimal}. These works provide key structural and algorithmic insights, but they differ from our setting, which allows fully adversarial arrivals and costs, nonlinear losses revealed only after action selection, and arbitrary real-valued allocation from current inventory under carry-over state dynamics. 

\paragraph{Inventory Management}
Yang et al. study an online inventory problem with time-varying linear prices and demand constraints, where the algorithm chooses purchases to satisfy demand and storage dynamics \cite{yang2020online}. They propose an algorithm and establish optimal competitive-ratio guarantees under linear costs \cite{yang2020online}. In contrast, we allow general convex costs on state-action pairs and analyze regret rather than worst-case competitive ratio.

Hihat et al. cast multi-product inventory control as an online convex optimization problem with stateful dynamics (e.g., perishability) and general convex losses \cite{hihat2023inventory}. Their guarantees rely on a mild non-degeneracy condition on demand, whereas our model permits fully adversarial arrivals and cost functions without structural assumptions.

Other recent work applies online learning to more specialized inventory settings. Guo et al. use online mirror descent for a multiproduct lost-sales system with capacity constraints and censored demand feedback \cite{guo2024omd}. Lugosi et al. analyze repeated newsvendor learning with adversarial censored demand \cite{lugosi2024hardness}. These models provide important guarantees in their domains but differ from our setting, which allows arbitrary real-valued allocation from current inventory each round under general convex costs with carry-over state.



\paragraph{Inventory Control and Supply Chain Management}
A mathematically parallel framework from operations research is \emph{Inventory Control}, which has optimized warehouse and supply chain logistics for over a century. The core state evolution equation---$I_{t+1} = \max(0, I_t + x_t - d_t)$, where $I_t$ is inventory, $x_t$ is replenishment, and $d_t$ is demand---is isomorphic to resource allocation under bounded storage, particularly when inflows and outflows are stochastic. Foundational work by Arrow, Karlin, and Scarf \cite{arrow1958studies} and Scarf's seminal result on optimal $(s,S)$ policies \cite{scarf1960optimality} established that dynamic programming and threshold-based rules characterize optimal allocation under demand and supply uncertainty. Modern treatments \cite{zipkin2000foundations,porteus2002foundations} systematize these results under holding costs and shortage penalties, paralleling the cost structures in energy and resource systems. While classical inventory optimization assumes knowledge of demand distributions (violating our adversarial assumption), the structural insights---especially the efficacy of threshold policies and the interpretation of capacity as a bounded state constraint---provide valuable intuition for both offline planning and extensions of our framework to partially stochastic settings.


\paragraph{Online Inventory }

State-constrained resource allocation is an active research area spanning online control, competitive analysis, and inventory-constrained optimization. From a \emph{competitive-analysis} perspective, prior work includes storage-management guarantees under worst-case costs \cite{Chau2016} and online linear optimization with inventory constraints \cite{Yang2020}. Broader inventory-constrained online optimization is developed in \cite{yang2020online,lin2019competitive,chau2016cost,neely-inventory-control-cdc2010}. For instance, \cite{yang2020online} treats linear objectives known at each round, whereas our setting permits nonlinear losses revealed only after acting. Related algorithmic approaches include one-way trading and online knapsack formulations \cite{cao2020optimal,el2001optimal}, which provide complementary insights for dynamic allocation.




\paragraph{Dam Theory and Storage Processes}
Classical dam theory provides an early stochastic framework for storage dynamics with random inflows and controlled releases under state-availability constraints. Representative references include classical storage-process and monograph treatments~\cite{moran1954dams,prabhu1998stochastic}. Relative to our adversarial online-control setting, this literature typically assumes stochastic input models and focuses on long-run or stationary performance criteria, rather than regret against a hindsight comparator. In our context, it serves primarily as historical and structural motivation for state-coupled allocation under feasibility constraints.


\paragraph{Online control}

A recent line of work in online learning studies \emph{online control}, where a dynamical system is operated under adversarial disturbances and costs. The objective is low policy regret relative to the best policy in hindsight, rather than classical stabilization or worst-case guarantees. Foundational results for linear systems and adversarial settings include \cite{agarwal2019online, agarwal2019logarithmic, hazan2020nonstochasticBandit, simchowitz2020improper}; see \cite{hazan_introduction_2025} for a modern overview.

A central idea in this literature is DAC, which parameterizes control actions as affine or linear functions of past disturbances. This yields tractable surrogates for policy optimization and enables efficient no-regret updates in adversarial environments. Extensions include partial observability \cite{lale2020logarithmic} and bandit feedback \cite{hazan2020nonstochasticBandit}. SDAC adapts this template to state-constrained resource systems, where feasibility is coupled over time through the dynamics.


From a \emph{control} perspective, Lyapunov-style methods and drift-based techniques are classical in this space. Representative results include near-optimal scheduling with finite-capacity states \cite{Huang2013} and online gradient/drift-plus-penalty hybrids with explicit utility-capacity tradeoffs \cite{Yu2019}. These results, together with the competitive-analysis line, emphasize the core difficulty: temporal coupling induced by constrained state evolution. Our model addresses that coupling in an adversarial setting while preserving feasibility.




\paragraph{Research Gap and Paper Scope} Adversarial online control provides a strong foundation for our problem class: it accommodates arbitrary time-varying disturbances and unknown convex costs, with regret guarantees against hindsight comparators and tractable online updates. DAC is central in this literature because it turns policy search into a tractable online convex optimization problem. However, for budget-coupled allocation dynamics, vanilla DAC does not enforce per-round feasibility by construction. Our framework fills this gap by replacing raw DAC with SDAC, a simplex-constrained DAC-style family that preserves feasibility defined by \eqref{eq:model-feasibility} while retaining adversarial online-learning tractability. We benchmark against the best fixed SDAC policy chosen in hindsight.




\subsection{Contributions}
The contributions are organized in two layers:
\begin{itemize}
    \item \textbf{Modeling and policy-class contribution.} We formulate adversarial online dynamic resource allocation with state-coupled feasibility and state-dependent costs, where the cost depends on the stored amount in the storage, and we introduce \emph{Simplex Disturbance-Action Control} (SDAC), a DAC-inspired simplex-constrained policy class to compete with. The SDAC \emph{enforces feasibility} by construction while preserving DAC-style expressiveness.
    \item \textbf{Algorithmic and analytical contribution.} We develop an exponentiated-gradient online controller over SDAC weights in which a clipping step \emph{enforces feasibility} of the online action, and we analyze regret against the best fixed SDAC comparator in hindsight, achieving sublinear regret with logarithmic dependence on policy dimension.
\end{itemize}

\subsection{Paper Organization}
Section~\ref{sec:model-setup} formalizes the adversarial allocation model, assumptions, and regret benchmark. Section~\ref{sec:sdac} introduces the SDAC policy class and summarizes its structural properties. Section~\ref{sec:online-algorithm} presents the online controller, Section~\ref{sec:regret-analysis} develops regret guarantees, and Section~\ref{sec:experiments} reports experiments.


\section{Notation}
\label{sec:notation}

We denote by $t\in\{1,2,\ldots,T\}$ the time index (total horizon $T$). The system state $b_t\in\mathbb{R}_{\ge 0}$ is the available resource, and the action $u_t\in\mathbb{R}_{\ge 0}$ satisfies $0 \le u_t \le b_t$. The exogenous inflow is $a_t\in[0,1]$, and the cost function is $c_t:\mathbb{R}_{\ge 0}^2\to\mathbb{R}$. The online algorithm is $\mathcal{A}$, and $\pi$ (or $\pi(w)$ when parameterized by $w$) denotes a policy from the comparator class $\Pi = \mathrm{SDAC}_H^{\kappa}$. The parameters $\kappa\in(0,1]$ (gain) and $H\in\mathbb{N}$ (memory horizon) are fixed design parameters. The weight vector $w$ lies in the simplex $\Delta_H = \{w\in\mathbb{R}_{\ge 0}^H : \sum_i w^{(i)}=1\}$. The Lipschitz constants for the cost function are $L_b,L_u>0$. The inner product is denoted $\langle x,y\rangle = x^T y$.

\section{Model and Problem Setup}
\label{sec:model-setup}

We consider a discrete-time allocation process over rounds $t\in\{1,2,\ldots,T\}$. The system state is denoted by $b_t\in\mathbb{R}_{\ge 0}$, where $b_t$ represents the available budget at the beginning of round $t$. The initial state is given as
\[
b_1 = 0.
\]

At each round, the controller chooses an allocation action $u_t\in\mathbb{R}_{\ge 0}$ subject to state feasibility:

\begin{equation}
\label{eq:model-feasibility}
0 \le u_t \le b_t.
\end{equation}
The system then receives an exogenous inflow $a_t\in[0,1]$ and evolves according to
\begin{equation}
\label{eq:model-dynamics}
b_{t+1} = b_t - u_t + a_t.
\end{equation}

After the action is selected, a convex cost function
\[
c_t: \mathbb{R}_{\ge 0}^2 \to \mathbb{R}
\]
is revealed, and the incurred cost is $c_t(b_t,u_t)$. We assume each $c_t$ is convex in $(b,u)$ and satisfies a Lipschitz condition on the relevant domain. There exist constants $L_b,L_u>0$ such that for all admissible $(b,u)$ and $(b',u')$,
\begin{equation}
\label{eq:model-lipschitz}
|c_t(b,u)-c_t(b',u')| \le L_b|b-b'| + L_u|u-u'|.
\end{equation}

The information pattern is adversarial and non-anticipatory. At time $t$, the controller chooses $u_t$ before observing $a_t$ and $c_t$. Hence both the inflow sequence $\{a_t\}_{t=1}^T$ and the cost sequence $\{c_t\}_{t=1}^T$ may be arbitrary.

\paragraph{Assumptions}
Equations \eqref{eq:model-feasibility}--\eqref{eq:model-lipschitz} define the model. For analysis, we collect the remaining standing assumptions as follows:
\begin{enumerate}
    \item \textbf{Inflow bounds:} $a_t\in[0,1]$ for all $t$.
    \item \textbf{Adversarial timing:} at round $t$, $u_t$ is chosen before observing $a_t$ and $c_t$.
    \item \textbf{Cost regularity:} each $c_t$ is convex in $(b,u)$ and obeys \eqref{eq:model-lipschitz}.
    \item \textbf{Feasible policy class:} both the online policy and comparator policies are required to satisfy \eqref{eq:model-feasibility} under dynamics \eqref{eq:model-dynamics}.
\end{enumerate}

For an online algorithm $\mathcal{A}$ producing trajectory $\{(b_t^{\mathcal{A}},u_t^{\mathcal{A}})\}_{t=1}^T$, define cumulative cost
\begin{equation}
\label{eq:def-cumulative-cost}
J_T(\mathcal{A}) \triangleq \sum_{t=1}^T c_t\bigl(b_t^{\mathcal{A}},u_t^{\mathcal{A}}\bigr).
\end{equation}
For fixed $(\kappa,H)$, let the comparator class be $\Pi\triangleq\mathrm{SDAC}_H^{\kappa}$ (defined in Section~\ref{sec:sdac}). For each $\pi\in\Pi$, let $\{(b_t^{\pi},u_t^{\pi})\}_{t=1}^T$ be the induced feasible trajectory under the same inflow sequence. Regret is defined by
\begin{equation}
\label{eq:def-regret}
\mathrm{Regret}(T)
\triangleq
\sum_{t=1}^T c_t\bigl(b_t^{\mathcal{A}},u_t^{\mathcal{A}}\bigr)
-
\min_{\pi\in\Pi}
\sum_{t=1}^T c_t\bigl(b_t^{\pi},u_t^{\pi}\bigr).
\end{equation}
While the online learner may update its SDAC weights over time, the benchmark in \eqref{eq:def-regret} is the best single fixed SDAC policy chosen in hindsight.
The objective is to design an online policy that preserves feasibility in \eqref{eq:model-feasibility} for all $t$ and achieves sublinear regret as defined in \eqref{eq:def-regret}.

\section{Simplex Disturbance-Action Control (SDAC)}
\label{sec:sdac}

\subsection{From Disturbance-Action Control to SDAC}
Disturbance-Action Control (DAC) parameterizes the action as a linear function of the current state and a finite window of past disturbances. This representation is central in adversarial online control because it accommodates arbitrary disturbance and cost sequences, yields robust regret guarantees, and enables efficient online convex-optimization updates. In our setting, a representative DAC form is
\begin{equation}
\label{eq:dac-param}
u_t = \kappa b_t + \sum_{i=1}^{H} m^{(i)} a_{t-i},
\end{equation}
where $\kappa\in\mathbb{R}$ is a gain, $H\in\mathbb{N}$ is the memory horizon, and $m\in\mathbb{R}^H$ is a disturbance filter.

For our constrained allocation model, raw DAC is insufficient: \eqref{eq:dac-param} does not generally enforce \eqref{eq:model-feasibility} at every round. In particular, the action can violate either bound, with $u_t>b_t$ or $u_t<0$. We therefore introduce \emph{Simplex Disturbance-Action Control} (SDAC), which replaces unconstrained DAC filters by a simplex-constrained family tailored to budget dynamics. SDAC is designed to keep the DAC advantages above (adversarial robustness, regret-based comparison, and tractable online updates) while adding feasibility by construction.


\subsection{SDAC Definition}
Fix $\kappa\in(0,1]$ and memory horizon $H\in\mathbb{N}$. Let
\begin{equation}
\label{eq:def-simplex}
\Delta_H \triangleq \left\{w\in\mathbb{R}_{\ge 0}^H:\; \sum_{i=1}^H w^{(i)}=1\right\}
\end{equation}
be the probability simplex. For any $w\in\Delta_H$, define the policy $\pi(w)$ by
\begin{equation}
\label{eq:def-sdac-policy}
u_t^{\pi(w)} \;\triangleq\; \kappa b_t^{\pi(w)} + \sum_{i=1}^{H} w^{(i)}(1-\kappa)^i a_{t-i},
\qquad a_\tau\equiv 0\;\text{for}\;\tau\le 0.
\end{equation}
The induced state trajectory satisfies
\begin{equation}
\label{eq:sdac-state-update}
b_{t+1}^{\pi(w)} = b_t^{\pi(w)} - u_t^{\pi(w)} + a_t,
\qquad b_1^{\pi(w)}=0.
\end{equation}

The resulting class is
\[
\mathrm{SDAC}_H^{\kappa}\triangleq\{\pi(w):w\in\Delta_H\},
\]
with $\Delta_H$ defined in \eqref{eq:def-simplex}.

First, feasibility is preserved within the class: for any fixed $w\in\Delta_H$, the induced trajectory under \eqref{eq:def-sdac-policy}--\eqref{eq:sdac-state-update} satisfies $0\le u_t^{\pi(w)}\le b_t^{\pi(w)}$ for all $t$ (see Lemma~\ref{lem:appendix-feasibility} in Appendix~\ref{sec:appendix-proofs}).

Second, SDAC is compatible with online convex optimization. For fixed inflows, both $b_t^{\pi(w)}$ and $u_t^{\pi(w)}$ are linear in $w$. Therefore, since each $c_t$ is convex in $(b,u)$, the induced cost is convex in $w$ (see Lemma~\ref{lem:appendix-surrogate} in Appendix~\ref{sec:appendix-proofs}). Regret minimization against the best fixed SDAC comparator thus becomes online convex optimization over the simplex.





\section{Online Algorithm}
\label{sec:online-algorithm}

We now define an online controller over SDAC weights using exponentiated-gradient updates on the simplex.

For each round $t$, let
\[
\widetilde{a}_t \triangleq \big((1-\kappa)^1 a_{t-1},\ldots,(1-\kappa)^H a_{t-H}\big)^\top \in \mathbb{R}^H,
\]
and let $w_t\in\Delta_H$ be the online weight vector.

\paragraph{Action selection}
The controller first forms the unconstrained SDAC-style action
\begin{equation}
\label{eq:alg-unconstrained-action}
\widehat{u}_t \triangleq \kappa b_t^{\mathcal A} + \langle w_t,\widetilde{a}_t\rangle,
\end{equation}
then enforces feasibility by clipping:
\begin{equation}
\label{eq:alg-feasible-action}
u_t^{\mathcal A} \triangleq \min\{\widehat{u}_t,\,b_t^{\mathcal A}\}.
\end{equation}
The state update is
\begin{equation}
\label{eq:alg-state-update}
b_{t+1}^{\mathcal A} = b_t^{\mathcal A} - u_t^{\mathcal A} + a_t,
\qquad b_1^{\mathcal A}=0.
\end{equation}

\paragraph{Feasibility of the online trajectory}
We verify that the algorithm satisfies the state-action constraints at every round. The key mechanism is the clipping step in \eqref{eq:alg-feasible-action}, which enforces per-round feasibility by truncating the unconstrained proposal to the currently available state. Since $w_t\in\Delta_H$, $\widetilde{a}_t\ge 0$ componentwise, and $\kappa\in(0,1]$, equation \eqref{eq:alg-unconstrained-action} gives $\widehat u_t\ge 0$. By \eqref{eq:alg-feasible-action},
\[
u_t^{\mathcal A}=\min\{\widehat u_t,b_t^{\mathcal A}\}\in[0,b_t^{\mathcal A}],
\]
so the action feasibility constraint holds whenever $b_t^{\mathcal A}\ge 0$. It remains to prove $b_t^{\mathcal A}\ge 0$ for all $t$. The base case is $b_1^{\mathcal A}=0$. If $b_t^{\mathcal A}\ge 0$, then $u_t^{\mathcal A}\le b_t^{\mathcal A}$ and $a_t\in[0,1]$, so from \eqref{eq:alg-state-update},
\[
b_{t+1}^{\mathcal A}=b_t^{\mathcal A}-u_t^{\mathcal A}+a_t\ge 0.
\]
Thus, by induction, $b_t^{\mathcal A}\ge 0$ and $0\le u_t^{\mathcal A}\le b_t^{\mathcal A}$ for all $t$.



\paragraph{Surrogate cost and gradient step}
Define the surrogate cost at round $t$ by
\begin{equation}
\label{eq:alg-surrogate-cost}
\ell_t(w) \triangleq c_t\bigl(b_t^{\pi(w)},u_t^{\pi(w)}\bigr), \qquad w\in\Delta_H,
\end{equation}
which is convex and Lipschitz by Lemma~\ref{lem:appendix-surrogate}. Let $g_t\in\partial \ell_t(w_t)$ be a subgradient.

The exponentiated-gradient update is
\begin{equation}
\label{eq:alg-eg-unnormalized}
\widetilde{w}_{t+1}^{(i)} \triangleq w_t^{(i)}\exp\!\bigl(-\eta g_t^{(i)}\bigr), \qquad i=1,\ldots,H,
\end{equation}
followed by simplex normalization:
\begin{equation}
\label{eq:alg-eg-normalized}
w_{t+1}^{(i)} \triangleq \frac{\widetilde{w}_{t+1}^{(i)}}{\sum_{j=1}^H \widetilde{w}_{t+1}^{(j)}}, \qquad i=1,\ldots,H.
\end{equation}

This yields a feasible online trajectory by \eqref{eq:alg-feasible-action} and performs online convex optimization over $\Delta_H$ through \eqref{eq:alg-surrogate-cost}--\eqref{eq:alg-eg-normalized}.


\section{Regret Analysis}
\label{sec:regret-analysis}


\section{Experiments}
\label{sec:experiments}

% TODO: EXPERIMENTENT BERTTER BE IN ENERGY HARVESTING FRAMEWORK

\appendix
\section{Appendix: Deferred Proofs}
\label{sec:appendix-proofs}

This appendix collects the deferred proofs for structural SDAC properties used in Section~\ref{sec:sdac}.

\begin{lemma}[State decomposition under SDAC]
\label{lem:appendix-state-decomp}
Fix $\kappa\in(0,1]$, $H\in\mathbb{N}$, and $w\in\Delta_H$. Let $\pi(w)$ be defined by \eqref{eq:def-sdac-policy} and \eqref{eq:sdac-state-update}. Then for all $t\ge 2$,
\begin{equation}
\label{eq:appendix-state-decomp}
b_t^{\pi(w)}
=
\sum_{i=1}^H w^{(i)}\!\left(\sum_{j=1}^{i}(1-\kappa)^{j-1}a_{t-j}\right)
= \langle w,\psi_t\rangle,
\end{equation}
where
\[
\psi_t^{(i)} \triangleq \sum_{j=1}^{i}(1-\kappa)^{j-1}a_{t-j}, \qquad i=1,\ldots,H.
\]
In particular, $0\le \psi_t^{(i)}\le 1/\kappa$. Consequently, since $w\in\Delta_H$ and
$b_t^{\pi(w)}=\langle w,\psi_t\rangle$, we also have $b_t^{\pi(w)}\ge 0$ for all $t\ge 2$.
\end{lemma}

\begin{proof}
We use induction on $t$. For $t=2$, $b_1^{\pi(w)}=0$ and $a_{1-i}=0$ for all $i\ge 1$. By the policy definition \eqref{eq:def-sdac-policy},
\(u_1^{\pi(w)}=\kappa b_1^{\pi(w)}+\sum_{i=1}^H w^{(i)}(1-\kappa)^i a_{1-i}=0\). Then, using the state update \eqref{eq:sdac-state-update},
\(b_2^{\pi(w)}=b_1^{\pi(w)}-u_1^{\pi(w)}+a_1=a_1\), which is exactly \eqref{eq:appendix-state-decomp} at time $t=2$.

Assume the claim holds at time $t\ge 2$. Using \eqref{eq:def-sdac-policy} and \eqref{eq:sdac-state-update},
\[
b_{t+1}^{\pi(w)}
= (1-\kappa)b_t^{\pi(w)} - \sum_{i=1}^H w^{(i)}(1-\kappa)^i a_{t-i} + a_t.
\]
Substituting the induction hypothesis gives
\[
\begin{aligned}
b_{t+1}^{\pi(w)}
&= (1-\kappa)\sum_{i=1}^H w^{(i)}\!\left(\sum_{j=1}^{i}(1-\kappa)^{j-1}a_{t-j}\right)
    - \sum_{i=1}^H w^{(i)}(1-\kappa)^i a_{t-i} + a_t \\
&= \sum_{i=1}^H w^{(i)}\!\left(\sum_{j=1}^{i}(1-\kappa)^j a_{t-j}\right)
    - \sum_{i=1}^H w^{(i)}(1-\kappa)^i a_{t-i} + a_t \\
&= \sum_{i=1}^H w^{(i)}\!\left(\sum_{j=1}^{i-1}(1-\kappa)^j a_{t-j}\right) + a_t \\
&= \sum_{i=1}^H w^{(i)}\!\left(a_t + \sum_{j=1}^{i-1}(1-\kappa)^j a_{t-j}\right) \\
&= \sum_{i=1}^H w^{(i)}\!\left(\sum_{j=1}^{i}(1-\kappa)^{j-1}a_{t+1-j}\right),
\end{aligned}
\]
which is exactly \eqref{eq:appendix-state-decomp} at time $t+1$. The bound $\psi_t^{(i)}\le 1/\kappa$ follows from $a_\tau\in[0,1]$ and the geometric series.
\end{proof}

\begin{lemma}[Linearity of SDAC action in $w$]
\label{lem:appendix-action-linear}
Under the setup of Lemma~\ref{lem:appendix-state-decomp}, for each $t$,
\[
u_t^{\pi(w)} = \langle w,\phi_t\rangle,
\]
where
\[
\phi_t^{(i)} \triangleq \kappa\psi_t^{(i)} + (1-\kappa)^i a_{t-i}
= \sum_{j=1}^{i}(1-\kappa)^{j-1}a_{t-j}.
\]
Moreover, $0\le \phi_t^{(i)}\le 1$ for all $i$.
\end{lemma}

\begin{proof}
From \eqref{eq:def-sdac-policy},
\[
u_t^{\pi(w)} = \kappa b_t^{\pi(w)} + \sum_{i=1}^H w^{(i)}(1-\kappa)^i a_{t-i}.
\]
Applying Lemma~\ref{lem:appendix-state-decomp}, we substitute
\(b_t^{\pi(w)}=\langle w,\psi_t\rangle=\sum_{i=1}^H w^{(i)}\psi_t^{(i)}\), so
\[
\begin{aligned}
u_t^{\pi(w)}
&= \kappa\sum_{i=1}^H w^{(i)}\psi_t^{(i)} + \sum_{i=1}^H w^{(i)}(1-\kappa)^i a_{t-i} \\
&= \sum_{i=1}^H w^{(i)}\bigl(\kappa\psi_t^{(i)} + (1-\kappa)^i a_{t-i}\bigr) \\
&= \sum_{i=1}^H w^{(i)}\phi_t^{(i)}
= \langle w,\phi_t\rangle.
\end{aligned}
\]

The lower bound on $\phi_t^{(i)}$ is immediate from $a_{t-i}\ge 0$ and $\kappa\le 1$. For the upper bound, use $a_\tau\in[0,1]$:
\[
\phi_t^{(i)}
\le \kappa\sum_{j=1}^{i}(1-\kappa)^{j-1} + (1-\kappa)^i
= 1.
\]
\end{proof}

\begin{lemma}[Feasibility of SDAC policies]
\label{lem:appendix-feasibility}
For every fixed $w\in\Delta_H$, the trajectory induced by \eqref{eq:def-sdac-policy}--\eqref{eq:sdac-state-update} satisfies
\[
0\le u_t^{\pi(w)}\le b_t^{\pi(w)} \quad \text{for all } t\ge 1.
\]
\end{lemma}

\begin{proof}
Nonnegativity of $u_t^{\pi(w)}$ follows from \eqref{eq:def-sdac-policy}, because each term in the sum is nonnegative: $\kappa b_t^{\pi(w)}\ge 0$ since $\kappa>0$ and $b_t^{\pi(w)}\ge 0$ (by Lemma~\ref{lem:appendix-state-decomp} for $t\ge 2$, and $b_1^{\pi(w)}=0$ by \eqref{eq:sdac-state-update}), while $w^{(i)}\ge 0$, $(1-\kappa)^i\ge 0$, and $a_{t-i}\ge 0$ for every $i$.

For the upper bound, Lemma~\ref{lem:appendix-state-decomp} implies
\[
(1-\kappa)b_t^{\pi(w)}
=
(1-\kappa)\sum_{i=1}^H w^{(i)}\!\left(\sum_{j=1}^{i}(1-\kappa)^{j-1}a_{t-j}\right)
=
\sum_{i=1}^H w^{(i)}\!\left(\sum_{j=1}^{i}(1-\kappa)^j a_{t-j}\right).
\]
Now rewrite the inner sum by collecting the coefficient of each $a_{t-i}$. Since all weights are nonnegative, for every $i$ we have
\[
\sum_{j=i}^{H} w^{(j)} \ge w^{(i)}.
\]
Therefore,
\[
(1-\kappa)b_t^{\pi(w)}
=
\sum_{i=1}^{H}\left(\sum_{j=i}^{H}w^{(j)}\right)(1-\kappa)^i a_{t-i}
\ge
\sum_{i=1}^{H}w^{(i)}(1-\kappa)^i a_{t-i},
\]
since $\sum_{j=i}^{H}w^{(j)}\ge w^{(i)}$ and $a_{t-i}\ge 0$. Therefore,
\[
u_t^{\pi(w)}
=
\kappa b_t^{\pi(w)} + \sum_{i=1}^{H}w^{(i)}(1-\kappa)^i a_{t-i}
\le
\kappa b_t^{\pi(w)} + (1-\kappa)b_t^{\pi(w)}
=
b_t^{\pi(w)}.
\]
\end{proof}

\begin{lemma}[Convexity and Lipschitz continuity of surrogate cost]
\label{lem:appendix-surrogate}
Define
\[
\ell_t(w) \triangleq c_t\bigl(b_t^{\pi(w)},u_t^{\pi(w)}\bigr), \qquad w\in\Delta_H.
\]
Then $\ell_t$ is convex on $\Delta_H$, and for all $w,w'\in\Delta_H$,
\[
|\ell_t(w)-\ell_t(w')|
\le
\left(\frac{L_b}{\kappa}+L_u\right)\|w-w'\|_1.
\]
\end{lemma}

\begin{proof}
By Lemma~\ref{lem:appendix-state-decomp} and Lemma~\ref{lem:appendix-action-linear},
\[
\ell_t(w)=c_t\bigl(\langle w,\psi_t\rangle,\langle w,\phi_t\rangle\bigr).
\]
The map $w\mapsto (\langle w,\psi_t\rangle,\langle w,\phi_t\rangle)$ is affine, and $c_t$ is convex, so $\ell_t$ is convex.

Using \eqref{eq:model-lipschitz},
\begin{align*}
|\ell_t(w)-\ell_t(w')|
&= \bigl|c_t\bigl(\langle w,\psi_t\rangle,\langle w,\phi_t\rangle\bigr)-c_t\bigl(\langle w',\psi_t\rangle,\langle w',\phi_t\rangle\bigr)\bigr| \\
&\le L_b\bigl|\langle w-w',\psi_t\rangle\bigr| + L_u\bigl|\langle w-w',\phi_t\rangle\bigr| \\
&\le L_b\|w-w'\|_1\|\psi_t\|_\infty + L_u\|w-w'\|_1\|\phi_t\|_\infty \\
&= \bigl(L_b\|\psi_t\|_\infty + L_u\|\phi_t\|_\infty\bigr)\|w-w'\|_1.
\end{align*}
Lemma~\ref{lem:appendix-state-decomp} gives $\|\psi_t\|_\infty\le 1/\kappa$, and Lemma~\ref{lem:appendix-action-linear} gives $\|\phi_t\|_\infty\le 1$, which yields the claim.
\end{proof}



\bibliographystyle{plain}
\bibliography{ref}

\end{document}
